% \subsection{RQ1: With What Precision Can Expected Behaviours Be Extracted from Source Code and Documentation?}
\subsection{RQ1 (Feasibility): With what precision can expected behaviours be extracted from source code and documentation?}
\label{sec:rq1}

% we should better frame this rq so the reviwers don't mix up that this is not part of the pipeline just to establish that behavioural gaps are valid gaps.
% If \toolname cannot reliably identify \emph{valid} behaviours given a method code and documentation, then any downstream analysis of behavioural gaps would be meaningless.
% Before quantifying behavioural gaps, we must first establish that the behaviours extracted by \toolname are valid. If the extracted behaviours do not accurately reflect expected behaviour from the code and documentation, any downstream analysis of behavioural gaps would be unreliable. Therefore, this RQ evaluates the precision of \toolname's behaviour extraction.

Before quantifying behavioural gaps, we must ensure that the behaviours the \toolname{} proof-of-concept surfaces are valid. 
%If the extracted behaviours do not accurately reflect expected behaviour from the code and documentation, any downstream analysis of behavioural gaps would be unreliable.
% Therefore, b
% Before quantifying the untested behaviours
To this end, we first investigate the precision of the \toolname{} behaviour extraction (Figure~\ref{fig:tool}(b)). 
% \toolname identified 20,729 candidate behaviours across all methods in our dataset. Manual inspection of all the identified behaviours would take a substantial amount of time and effort. We instead build an LLM judge to estimate validity at scale. We first validate the judge against human annotators, then use it to derive the extraction precision of \toolname.
\toolname{} identified 20,729 behaviours across all documented methods in our dataset.
% Manually inspecting all identified behaviours would take a substantial amount of time and effort. 
To evaluate whether these behaviours are correct, 
we adopt a two-step approach: we first establish human consensus labels to determine how consistently behaviours can be labelled from documentation and source code.
We next build an LLM-based judge, compare its labelling performance against the prior human labels, and apply this judge against all identified behaviours to assess the overall precision of \toolname{}.

% as the ground truth, then validate an LLM judge against them, and use the validated judge as a scaling instrument to estimate extraction precision across all identified behaviours.

% we build an LLM judge\footnote{This judge is not part of the original pipeline of \toolname} to classify behaviours at scale, validate the judge against human annotators, and then use the validated judge as a measurement instrument to estimate overall extraction precision.

\smallskip
\noindent\textbf{Behaviour assessment judge:} The judge is an agentic system built on the Ministral 3:14B reasoning model with its default configuration. We use a different model family (Ministral) as the judge from the one used for behaviour extraction (Qwen) to mitigate self-enhancement bias~\cite{zheng2023judging}.
% The judge serves solely as a scaling instrument. 
%Its classifications define strata for estimation, with per-stratum validity rates derived entirely from human consensus labels rather than from the judge itself. 
We choose this model for its strong reasoning performance relative to similar-sized models, which is important for behaviour validation~\cite{liu2026ministral}. 
Substituting the model with a more advanced reasoning model should only increase the judge's precision. 

The judge is provided a method behaviour, the method's documentation and implementation, and can query other method implementations as needed via tool use. 
% RTH: overlooks the 'testability' of the behaviour...
% A behaviour is defined in the base prompt as valid if it is inferable from the documentation and the implementation (documentation is prioritized), and expressed as testable over inputs, outputs, or exceptions. 
The judge then assesses whether the provided behaviour would be expected. 
Across all projects, the judge labelled 19,829 (95.7\%) behaviours as valid and 900 (4.3\%) as invalid.
Table~\ref{tab:method_behaviour_counts} reports the judge's output alongside \toolname{}'s extracted counts. % for each project. 
% In the base prompt, we define a valid behaviour as one that is supported by the method’s documentation and implementation, and is testable. 




\smallskip
\noindent\textbf{Validating behaviour judge against human labels:}
%
Recent work has shown that LLMs can effectively label software engineering tasks when properly validated~\cite{ahmed2025can, 10.1145/3797276, zhou2025llm}. We follow this practice and validate our judge through manual labelling before trusting it at scale.
% To validate the judge's performance, we conduct a manual evaluation. 
As the judge’s predictions are imbalanced (19,829 valid vs. 900 invalid), a proportional random sample would 
over-sample valid judgments. % compared to invalid judgments.
% yield only about 5 invalid cases, which is insufficient for reliable per-stratum estimation.
We therefore used a disproportional stratified random sample~\cite{rs11020185} of 100 behaviours, taking 50 from each stratum (valid and invalid) to ensure the minority class is sufficiently represented. %for estimating per-stratum agreement rates independently of stratum size. 
Within each stratum, we further distribute samples evenly across the 10 projects. % to reduce sample bias. 
% without altering the stratum sizes used in estimation.

% To maintain balanced coverage across each project, we also ensured each of the 10 projects contributed 10 samples.
% i am not confident framing it correctly here so commenting it out
% Although not based on a formal power analysis, this sample size is consistent with standard practice in empirical SE research, where extensive manual evaluation typically limits the proportion of artifacts analyzed while still supporting representative and reliable conclusions.~\cite{Muttakin2026TheSO}

Two authors, each with several years of Java experience, independently labelled each of the 100 behaviours as \textit{valid} or \textit{invalid} based on the method and its documentation. 
The two labellers reached \emph{substantial} agreement (Cohen's $\kappa=0.77$, Krippendorff's $\alpha=0.79$)~\cite{mchugh2012interrater}.
A post-hoc power analysis ($\alpha=0.05$, power=0.90) indicates that 43 samples are sufficient to detect this level of agreement~\cite{doi:10.1177/02537176261422290}, well within our sample of 100. 
The labellers resolved any disagreements collaboratively to produce a consensus label for each behaviour to serve as a ground truth against which to evaluate the judge.

% \noindent\textbf{How well does the judge agree with consensus?} 
% To measure this we compare the judge's original prediction with the consensus labels. This gives $\kappa = 0.68$, $\alpha = 0.68$, which represents \emph{substantial} agreement with the consensus labelled data. 
Comparing the judge's original prediction to these consensus labels yields $\kappa = 0.68$ and $\alpha = 0.68$, again representing \emph{substantial} agreement.
The post-hoc power analysis ($\alpha=0.05$, power=0.90) indicates that a minimum of 89 samples are required to detect this level of agreement.

%, and our sample of 100 exceeds this threshold. 
%The judge’s per‑project agreement with human consensus ranged from 85\% to 100\%, with an overall agreement of 96\%. This suggests the judge performs consistently across projects.
% and it showed that required minimum sample  size is 89 and our sample size already exceeds this. 

% RTH: this is good detail, but might be more than is needed:
% v2: Since the judge achieves substantial agreement with human consensus labels, we therefore treat the judge's classifications as reliable strata for estimation over the full behaviour set. Importantly, the judge's classifications define strata only. The per‑stratum validity rates $\hat{p}_v$ and $\hat{p}_i$ are derived directly from human consensus labels, thereby correcting for any judge misclassification.
% v1: the judge's classifications serve only as strata definitions, not as ground truth labels. The per-stratum validity rates $\hat{p}_v$ and $\hat{p}_i$ are derived directly from the human consensus labels and already correct for the judge's observed misclassification rates. 
%v2: Thus, even if the judge misclassifies behaviours across strata, the stratified estimator corrects for the per‑stratum validity rates derived from human consensus, yielding an estimate that reflects the human consensus labels.

% The set of expected behaviours is inherently subjective. Even two human annotators agreement didn't achieve perfect $\kappa$ initially and there are no complete ground truth for this task. Therefore recall is not measurable for this task and we instead on precision. A high precision means when \toolname identifies a behaviour we can be confident it's a valid one and represents a lower bound of the expected behaviours.

\begin{table}[h]
\centering
\caption{\toolname{} extracted behaviours and judged validity.}
\label{tab:method_behaviour_counts}
\begin{tabular}{lrrrr}
\toprule
\textbf{Project} & \textbf{\# Behaviours} & \textbf{\# Valid} & \textbf{\# Invalid (\%)} \\
\midrule
commons-cli         & 424 & 402 & 22 (5\%)\\
commons-codec       & 1,218 & 1,158 & 60 (5\%) \\
commons-compress    & 1,707 & 1,614 & 93 (5\%) \\
commons-jxpath      & 183 & 177 & 6 (3\%) \\
commons-collections & 2,816 & 2,681 & 135 (5\%) \\
commons-lang        & 7,719 & 7,380 & 339 (4\%)\\
commons-csv         & 361 & 348 & 13 (3\%) \\
gson                & 554 & 536 & 18 (3\%) \\
jfreechart          & 4,532 & 4,366 & 166 (4\%) \\
jsoup               & 1,215 & 1,167 & 48 (4\%) \\
\midrule
Total               & 20,729 & 19,829 & 900 (4\%) \\
\bottomrule
\end{tabular}
%\vspace{-1em}
\end{table}

\smallskip
\noindent\textbf{Estimating precision over all extracted behaviours:}
Since the judge serves as our measurement instrument over all behaviours ($AB=20,729$), its precision directly determines our confidence in \toolname{}'s extracted behaviours.
Recall cannot be calculated because expected behaviours lack a discrete ground truth; their decomposition into individual behaviours is inherently subjective. 
We therefore evaluate precision, as determining whether an extracted behaviour is valid constitutes a well-defined binary judgment.
A high precision implies that when \toolname identifies a behaviour, we can be confident it is valid and expected given a method and its documentation. 

We use the consensus labels from the 100 samples to estimate per-stratum validity rates.
Among the 50 samples drawn from the judge's valid stratum, 48 are confirmed valid compared to the human labels ($\hat{p}_{valid} = 0.96$). 
Among the 50 samples from the invalid stratum, 14 are confirmed valid ($\hat{p}_{invalid} = 0.28$). This reflects the judge's conservative tendency to flag some valid behaviours as invalid. 

Specifically, the judge tended to reject behaviours that were implied rather than explicitly stated in the documentation.
For example, flagging a behaviour about negative inputs as invalid when the Javadoc only discussed positive ones. 
The judge also struggled with ambiguous descriptions that human annotators could resolve from context. 
The former reflects strict adherence to explicit documentation, the latter conservative rejection of under-specified wording, both of which cause the estimated precision to be a lower bound on \toolname{}'s true precision. 
% As both failure modes cause the judge to under-count valid behaviours, the estimated precision for \toolname{} should be interpreted as a lower bound on its true precision.

%\rth{not clear whether the hats over AB/VB are actually needed.}

As the invalid strata comprises only 4.3\% of the judgments, the overall precision remains high, despite the judge being conservative with the invalid strata.
% Since both error patterns are systematic, the stratified estimator corrects for them directly.
% Let $N_v = 19,829$ be the size of the judge's valid stratum and $N_i = 900$ the size of the invalid stratum. 
Applying the stratified proportion estimator~\cite{scheaffer2012elementary}(Ch~5.6),
%with $N_v = 19{,}829$ and $N_i = 900$, 
the estimated total number of valid behaviours ($V\!B_{est}$) surfaced by \toolname{} %across the dataset 
is:
% chapter 5 eq 5.13 of the book
\[
V\!B_{est} = n_{valid} \cdot \hat{p}_{valid} + n_{invalid} \cdot \hat{p}_{invalid} \hspace{0.6em}% = 19{,}288 
\]
\vspace{-1.5em}
\[
V\!B_{est} = 19{,}829 \cdot 0.96 + 900 \cdot 0.28 = 19{,}288
\]
% \[
% = 19{,}829 \times 0.96 + 900 \times 0.28
% % = 19,\!036 + 252 
% = 19{,}288
% \]
Precision is then: 
\[
\hat{P_{beh}} = \frac{VB_{est}}{AB} = \frac{19,\!288}{20,\!729} \approx 0.931
\]
% chapter 5 eq 5.14 of the book
% \[
% \text{Var}(\hat{P}) = \frac{1}{N^2} \sum_{h \in \{v,i\}} 
% N_h^2 \cdot \frac{1 - f_h}{n_h} \cdot \hat{p}_h(1-\hat{p}_h)
% \]
% where $f_h = n_h / N_h$ is the sampling fraction for stratum $h$. Substituting the values yields $\text{SE}(\hat{P}) \approx 0.0266$, and the 95\% confidence interval is 
% $0.931 \pm 1.96 \times 0.0266 = [87.8\%, 98.3\%]$
%To quantify uncertainty around this precision estimate, we apply the 
% The 95\% confidence interval, computed using the
%stratified variance estimator with finite population correction: 
To bound precision uncertainty, we apply a stratified variance estimator with finite population correction:
\[
\text{Var}(\hat{P_{beh}}) = \frac{1}{N^2} \sum_{h \in \{v,i\}} 
N_h^2 \cdot \frac{1 - f_h}{n_h} \cdot \hat{p}_h(1-\hat{p}_h)
\] where $f_h = n_h / N_h$ is the sampling fraction for stratum $h$. 
%Substituting the values 
This yields standard error $\text{SE}(\hat{P_{beh}}) \approx 0.0266$.



\smallskip
\noindent\textbf{Implication:} This analysis shows that \toolname{} is able to reliably extract real behaviours from software methods and their documentation with high precision. 
The methodology also demonstrates how a carefully-developed and validated automated judge can be shown to have substantial inter-rater reliability compared to multiple human labellers, enabling analysis of whole datasets instead of small manual samples.

The judge's errors were not random. Failures in behaviour extraction concentrated in cases where behaviours were implied by context rather than explicitly stated, or where documentation was ambiguous. 
Both failure modes cause underestimation: implied behaviours that the judge rejects are real behaviours that do not enter the reference set; ambiguous behaviours flagged as invalid reduce the denominator of confirmed gaps. 
As new models are developed that could be used for the behaviour extractor or judge, more of these cases could be resolved correctly, increasing the estimated number of valid behaviours that can be checked against existing test cases.
% Our reported 93.1\% figure is therefore conservative with respect to model capability.


%Across all evaluated projects, \toolname{} demonstrates consistent performance, yielding a conservative overall behaviour detection precision of $93.1\%\pm2.7\%$.
% , and the 95\% confidence interval is $0.931 \pm 1.96 \times 0.0266 = [87.8\%, 98.3\%]$.
% This yields a 95\% confidence interval of $[87.8\%, 98.3\%]$.
% Taking the lower bound conservatively, at least $87.8\%$ of the behaviours identified by \toolname are valid with 95\% confidence.

% Since the set of expected behaviours is inherently subjective and no complete 
% ground truth exists for this task, recall is not directly measurable. We 
% therefore focus on precision. A high precision means that when \toolname 
% identifies a behaviour we can be confident it is valid. The estimated 
% precision of $93.1\%$ (conservatively $87.8\%$) thus represents a lower bound 
% on the true rate of valid behaviours identified by the tool.

% RTH: let's leave this out and just lean on the SE above
%\smallskip
%\noindent\textbf{Sensitivity to the small invalid stratum.} Because the invalid stratum consists of only $4.3\%$ of all behaviours, uncertainty in $\hat{p}_i$ has limited influence on the overall estimate. By comparison, uncertainty in $\hat{p}_v$ is negligible.\footnote{$\hat{p}_v(1-\hat{p}_v)/n_v = 0.96 \times 0.04 / 50 \approx 0.001$.} Varying $\hat{p}_i$ across its full range $[0\%,100\%]$ yields precision estimates between $91.8\%$ (at $\hat{p}_i = 0\%$) and $96.2\%$ (at $\hat{p}_i = 100\%$). Thus, even under the worst-case assumption that none of the invalid-stratum behaviours are truly valid, \toolname's precision remains above $91.8\%$.
% \[
% \hat{P}(p_i) = \frac{19,\!829 \times 0.96 + 900 \times p_i}{20,\!729}.
% \]

\begin{rqbox}{RQ1: Precision of \toolname{}'s behaviour extraction.}
%\textbf{Summary.} 
Across all evaluated projects, \toolname{} conservatively yields a behaviour detection precision of $93.1\%\pm2.7\%$. 
% For the analyzed dataset, this means at least 19,288 behaviours were identified across the 10 projects.
This high precision and low variance demonstrate that \toolname{} can consistently extract valid expected behaviours, establishing a reliable foundation for downstream test adequacy analysis.
%Using the consensus-validated judge as a measurement instrument, we estimate that \toolname extracts expected behaviours with a precision of 93.1\%.
% (95\% CI: [87.8\%, 98.3\%]. 
%Varying $\hat{p}_i$ across its full range [0, 1] shifts the overall precision only between 91.8\% and 96.2\%, since the invalid stratum represents just 4.3\% of all extracted behaviours (900 of 20,729). We therefore conservatively report that at least 91.8\% of the behaviours identified by \toolname are valid, giving us confidence that the behavioural gaps reported in RQ2 reflect genuine untested behaviours, not spurious or imaginary ones.
\end{rqbox}

% \rth{not sure whether the widehat's are necessary for AB/TB above. also, let's look at the human consensus precision and decide if we want to explicitly mention this to contextualize the precision above.}